\section*{Введение}
\addcontentsline{toc}{section}{Введение}

Современные организации активно используют облачные файловые хранилища для совместной
работы с документами. Seafile~— открытая (open-source) платформа для хранения и
синхронизации файлов, которую можно развернуть на собственном сервере~\cite{seafile_docs}.

Целью данной практической работы является установка и настройка сервера Seafile~\cfgSeafileVersion{}
на базе Ubuntu~22.04 в качестве виртуальной машины \texttt{seafile}
внутри локальной сети \texttt{192.168.\cfgLabStudentN.0/24}.

В рамках работы решаются следующие задачи:
\begin{itemize}
  \item настройка статического IP-адреса (\texttt{\cfgSeafileIP}) через Netplan;
  \item регистрация ВМ в домене \texttt{\cfgStudentDomain} (DNS-записи);
  \item установка MariaDB и необходимых Python-зависимостей;
  \item установка и инициализация сервера Seafile;
  \item настройка обратного прокси Nginx~\cite{nginx_docs};
  \item настройка автозапуска служб через systemd;
  \item установка клиента Seafile на Desktop и синхронизация файлов.
\end{itemize}

Стенд: виртуальная машина \texttt{seafile} (Ubuntu~22.04, VirtualBox,
адаптер <<Внутренняя сеть>>), клиентская ВМ \texttt{Desktop1} (Ubuntu Desktop).
